\chapter{Regression Discontinuity Design}
\label{cap:rdd}
% ── index ───────────────────────────────────────────────────────
\index{regression discontinuity|textbf}
\index{RDD|see{regression discontinuity}}
\index{rd@\texttt{rd()}|textbf}
\index{cutoff threshold}
\index{running variable}

% ─────────────────────────────────────────────────────────────────────────────
\section{The model}

Regression Discontinuity Design (RDD) identifies causal effects when
treatment is assigned mechanically based on a threshold of a continuous
observable variable --- the \textbf{running variable} $x$:

\[
  D_i = \mathbf{1}[x_i \geq c]
\]

The estimator exploits the fact that individuals immediately below and above the
threshold $c$ are comparable in all respects except the treatment received. The
identified effect is the \textbf{LATE at the cutoff} (\emph{Local Average Treatment
Effect at the cutoff}):

\[
  \tau_{\mathrm{RDD}}
  = \lim_{x \downarrow c} E[y \mid x]
  - \lim_{x \uparrow c} E[y \mid x]
\]

\begin{nota}
  The \emph{Sharp} RDD (above) assumes that $D$ changes deterministically at $c$.
  The \emph{Fuzzy} RDD generalizes to when $D$ increases
  \emph{discontinuously} at $c$, without a jump from 0 to 1 --- it requires a
  first stage and is equivalent to a local IV.
\end{nota}

% ─────────────────────────────────────────────────────────────────────────────
\section{Identification}

The central assumption is \textbf{continuity of potential outcomes} at
$c$:

\[
  E[Y^{(0)} \mid x] \quad \text{and} \quad E[Y^{(1)} \mid x]
  \quad \text{are continuous at } x = c
\]

This guarantees that any discontinuity in $E[y \mid x]$ at point $c$ is
caused by the jump in $D$, and not by another factor that also changes at $c$.

The assumption is violated if agents \emph{manipulate} $x$ to place themselves above or
below the threshold --- students who revise answers to cross the passing cutoff,
for example. The McCrary density test (continuity test for the density of $x$ at $c$)
checks for suspicious bunching of observations on one side of the threshold.

% ─────────────────────────────────────────────────────────────────────────────
\section{The estimator}

The standard estimator is \textbf{local linear regression} fit separately on each
side of the threshold, within a \emph{bandwidth} $h$:

\[
  \hat\tau = \hat\mu_+ - \hat\mu_-
\]

where $\hat\mu_+$ and $\hat\mu_-$ are the intercepts of the local linear
regressions on the right side ($x \in [c, c+h)$) and left side ($x \in (c-h, c)$),
respectively, weighted by the triangular kernel:

\[
  K\!\left(\frac{x - c}{h}\right) = \left(1 - \frac{|x-c|}{h}\right)
  \mathbf{1}[|x-c| \leq h]
\]

The \emph{bandwidth} $h$ controls the bias-variance trade-off: wider windows reduce
variance but increase bias (local linearity may be a poor approximation);
narrower windows are more local but have fewer observations. Hayashi
selects $h$ automatically by the MSE-optimal criterion (Imbens-Kalyanaraman)
when \lstinline{bw} is not specified.

% ─────────────────────────────────────────────────────────────────────────────
\section{DGP Verification}

\subsection*{Exact recovery (no noise)}

The DGP uses a running variable $x \sim \mathrm{Uniform}(-1, 1)$, threshold at
$c = 0$ and true effect $\tau = 2{,}0$. The outcome function is linear
in $x$ on both sides of the threshold, without noise, so RDD should recover
$\tau$ exactly.

\[
  y_i = 1{,}0 + 0{,}5\,x_i + 2{,}0\,D_i, \qquad D_i = \mathbf{1}[x_i \geq 0]
\]

\begin{lstlisting}[caption={Noise-free RDD DGP --- exact recovery}]
seed(42)
input df
id
1
end
for i in 1..=10 { df = append(df, df) }
df |> filter(_n <= 1000)
generate df x = runiform() * 2 - 1
generate df d = x >= 0
generate df y = 1.0 + 0.5*x + 2.0*d
rd(y ~ x, 0.0, df)
\end{lstlisting}

\begin{lstlisting}[language={}, basicstyle=\ttfamily\small,
  frame=none, numbers=none, backgroundcolor=\color{codbg},
  xleftmargin=0pt, xrightmargin=0pt, breaklines=false]
══════════════════════════════════════════════════════════════════════
 Regression Discontinuity  —  Sharp  —  Local Linear (p=1)
══════════════════════════════════════════════════════════════════════
 Outcome: y                   Running var: x
 Cutoff: 0.0000   Bandwidth: 0.0571   Kernel: Triangular
 Obs total: 72   N left: 37   N right: 35
──────────────────────────────────────────────────────────────────────
 Treatment Effect (τ̂):
       2.0000   SE     0.0000   z 2.443e15   P>|z|   0.0000  ***
 95% CI: [2.0000, 2.0000]
══════════════════════════════════════════════════════════════════════
 *** p<0.01  ** p<0.05  * p<0.10
\end{lstlisting}

The estimator recovers $\hat\tau = 2{,}0000$ exactly. The bandwidth of
$0{,}0571$ was selected automatically --- of the 1000 observations, 72 lie
within the window around the cutoff.

\subsection*{With idiosyncratic noise}

Adding idiosyncratic $\varepsilon$, the LATE remains identified:

\begin{lstlisting}[caption={RDD DGP with noise}]
generate df e = rnormal()
generate df y = 1.0 + 0.5*x + 2.0*d + e
rd(y ~ x, 0.0, df)
\end{lstlisting}

\begin{lstlisting}[language={}, basicstyle=\ttfamily\small,
  frame=none, numbers=none, backgroundcolor=\color{codbg},
  xleftmargin=0pt, xrightmargin=0pt, breaklines=false]
══════════════════════════════════════════════════════════════════════
 Regression Discontinuity  —  Sharp  —  Local Linear (p=1)
══════════════════════════════════════════════════════════════════════
 Outcome: y                   Running var: x
 Cutoff: 0.0000   Bandwidth: 0.0741   Kernel: Triangular
 Obs total: 94   N left: 43   N right: 51
──────────────────────────────────────────────────────────────────────
 Treatment Effect (τ̂):
       2.0116   SE     0.1191   z   16.893   P>|z|   0.0000  ***
 95% CI: [1.7782, 2.2449]
══════════════════════════════════════════════════════════════════════
 *** p<0.01  ** p<0.05  * p<0.10
\end{lstlisting}

With noise, $\hat\tau = 2{,}0116$ --- the deviation from $2{,}0$ is sampling
variance. The 95\% CI $[1{,}78;\; 2{,}24]$ covers the true $\tau = 2{,}0$.
The bandwidth was $0{,}0741$: with more variance, the MSE-optimal criterion
prefers a narrower window to control bias.

% ─────────────────────────────────────────────────────────────────────────────
\section{Syntax}

\begin{lstlisting}
rd(y ~ x, cutoff, df)
rd(y ~ x, cutoff, df, bw=0.3)
rd(y ~ x, cutoff, df, bw=0.3, poly=2)
\end{lstlisting}

\begin{center}
\begin{tabular}{lll}
\toprule
\textbf{Parameter} & \textbf{Default} & \textbf{Description} \\
\midrule
\texttt{cutoff}  & required    & threshold value in $x$ \\
\texttt{bw}      & auto (IK)   & window width $h$ \\
\texttt{poly}    & \texttt{1}  & degree of local polynomial \\
\texttt{kernel}  & \texttt{triangular} & \texttt{uniform}, \texttt{epanechnikov} \\
\bottomrule
\end{tabular}
\end{center}

\begin{lstlisting}[caption={Typical workflow --- test score cutoff and scholarship}]
load "school.csv" as df

// score centered at threshold 60
generate df xc = score - 60

let m = rd(approved ~ xc, 0.0, df)
display m
\end{lstlisting}

\section{Capturing the result}

\begin{lstlisting}
let m = rd(y ~ x, 0.0, df)
display m
\end{lstlisting}

\begin{nota}
  In real applications, validate the RDD with falsification tests:
  (1)~estimate the model using \emph{pre-treatment covariates} as
  the outcome --- there should be no discontinuity at $c$;
  (2)~test discontinuities at \emph{placebo thresholds} ($c \pm \delta$)
  --- they should not be significant; (3)~check the density of $x$ at
  $c$ (McCrary) to detect manipulation of the running variable.
\end{nota}

\section{Empirical validation}

The DGP in this chapter with a uniform running variable and cutoff at 0 is
used in the \texttt{rdd\_book} validation case in \texttt{validation/cases/}.
The case estimates the treatment effect by RDD in \hayashi{} and compares it
with R and Python reference implementations. Run \lstinline{hay validate} to
see the current status.
